\chapter{Casos de uso - Primera iteración}
\label{chap:anexo_casos_uso_it1}

Este anexo presenta la especificación detallada de los casos de uso desarrollados durante la primera iteración del sistema, incluyendo flujos principales, alternativos y de excepción, junto con precondiciones y postcondiciones específicas.

\section{CU-01 - Captura básica de eventos de entrada}
\label{sec:cu01_captura_eventos}

\begin{table}[H]
\centering
\resizebox{\textwidth}{!}{
\rowcolors{2}{white}{udcgray!25}
\begin{tabular}{|p{3cm}|p{12cm}|}
\hline
\rowcolor{udcpink!25}
\textbf{Campo} & \textbf{Descripción} \\
\hline
\textbf{Identificador} & CU-01 \\
\hline
\textbf{Nombre} & Captura básica de eventos de entrada para correlación posterior \\
\hline
\textbf{Actores} & Investigador BCI (principal), Usuario del sistema (secundario) \\
\hline
\textbf{Descripción} & El sistema debe capturar eventos de teclado y ratón con precisión temporal para facilitar correlación posterior con señales cerebrales \\
\hline
\textbf{Precondiciones} & 
\begin{minipage}[t]{11cm}
\begin{itemize}
    \item El sistema está iniciado correctamente
    \item Los dispositivos de entrada están conectados y funcionando
    \item Se ha configurado el filtro de teclas deseado
    \item El modo de tracking de ratón está seleccionado
\end{itemize}
\end{minipage} \\
\hline
\textbf{Postcondiciones} & 
\begin{minipage}[t]{11cm}
\begin{itemize}
    \item Los eventos están registrados con timestamps precisos
    \item Los datos están disponibles para exportación
    \item Las visualizaciones muestran actividad en tiempo real
    \item El sistema mantiene un histórico de eventos capturados
\end{itemize}
\end{minipage} \\
\hline
\end{tabular}
}
\caption{Especificación del caso de uso CU-01}
\label{tab:cu01_especificacion}
\end{table}

\subsection{Flujo principal}
\label{subsec:cu01_flujo_principal}

\begin{table}[H]
\centering
\resizebox{\textwidth}{!}{
\rowcolors{2}{white}{udcgray!25}
\begin{tabular}{|c|p{13cm}|}
\hline
\rowcolor{udcpink!25}
\textbf{Paso} & \textbf{Acción} \\
\hline
1 & El investigador BCI inicia la aplicación de captura de entrada \\
\hline
2 & El sistema inicializa los listeners de teclado y ratón mediante pynput \\
\hline
3 & El investigador configura el filtro de teclas específicas a capturar \\
\hline
4 & El investigador selecciona el modo de tracking de ratón (axis/angle/delta/joystick) \\
\hline
5 & El investigador inicia la captura de eventos \\
\hline
6 & El sistema comienza a capturar eventos con timestamps de alta precisión \\
\hline
7 & Los eventos se procesan y distribuyen a visualizadores en tiempo real \\
\hline
8 & El usuario del sistema realiza acciones que generan eventos \\
\hline
9 & El sistema registra cada evento con timestamp y metadatos asociados \\
\hline
10 & Los visualizadores actualizan la interfaz mostrando actividad \\
\hline
11 & El investigador detiene la captura cuando considera apropiado \\
\hline
12 & El sistema finaliza la captura y prepara datos para exportación \\
\hline
\end{tabular}
}
\caption{Flujo principal del caso de uso CU-01}
\label{tab:cu01_flujo_principal}
\end{table}

\subsection{Flujos alternativos}
\label{subsec:cu01_flujos_alternativos}

\begin{table}[H]
\centering
\resizebox{\textwidth}{!}{
\rowcolors{2}{white}{udcgray!25}
\begin{tabular}{|p{3cm}|p{11cm}|}
\hline
\rowcolor{udcpink!25}
\textbf{Flujo} & \textbf{Descripción} \\
\hline
\textbf{FA-01: Cambio configuración durante captura} & 
\begin{minipage}[t]{10.5cm}
\begin{itemize}
    \item En el paso 6, el investigador modifica filtros de teclas
    \item El sistema aplica cambios sin interrumpir la captura
    \item Se registra el cambio con timestamp para trazabilidad
    \item Continúa en paso 7
\end{itemize}
\end{minipage} \\
\hline
\textbf{FA-02: Cambio modo tracking} & 
\begin{minipage}[t]{10.5cm}
\begin{itemize}
    \item En el paso 6, el investigador cambia modo de tracking de ratón
    \item El sistema actualiza el procesamiento de movimiento
    \item Los visualizadores adaptan representación al nuevo modo
    \item Continúa en paso 7
\end{itemize}
\end{minipage} \\
\hline
\textbf{FA-03: Pausa temporal} & 
\begin{minipage}[t]{10.5cm}
\begin{itemize}
    \item En cualquier momento durante captura, investigador pausa
    \item El sistema suspende captura manteniendo configuración
    \item Al reanudar, continúa con timestamp actualizado
    \item Mantiene coherencia en secuencia de eventos
\end{itemize}
\end{minipage} \\
\hline
\end{tabular}
}
\caption{Flujos alternativos del caso de uso CU-01}
\label{tab:cu01_flujos_alternativos}
\end{table}

\subsection{Flujos de excepción}
\label{subsec:cu01_flujos_excepcion}

\begin{table}[H]
\centering
\resizebox{\textwidth}{!}{
\rowcolors{2}{white}{udcgray!25}
\begin{tabular}{|p{3cm}|p{11cm}|}
\hline
\rowcolor{udcpink!25}
\textbf{Excepción} & \textbf{Manejo} \\
\hline
\textbf{FE-01: Fallo dispositivo entrada} & 
\begin{minipage}[t]{10.5cm}
\begin{itemize}
    \item El sistema detecta desconexión de dispositivo
    \item Muestra alerta al investigador indicando dispositivo afectado
    \item Continúa captura con dispositivos disponibles
    \item Registra evento de fallo para análisis posterior
\end{itemize}
\end{minipage} \\
\hline
\textbf{FE-02: Error permisos sistema} & 
\begin{minipage}[t]{10.5cm}
\begin{itemize}
    \item pynput no puede acceder a eventos del sistema
    \item El sistema muestra mensaje explicativo sobre permisos
    \item Proporciona instrucciones específicas del sistema operativo
    \item Termina ejecución de forma controlada
\end{itemize}
\end{minipage} \\
\hline
\textbf{FE-03: Saturación buffer} & 
\begin{minipage}[t]{10.5cm}
\begin{itemize}
    \item Eventos llegan más rápido que procesamiento
    \item Sistema activa descarte selectivo de eventos duplicados
    \item Mantiene eventos únicos y registra pérdidas
    \item Ajusta automáticamente frecuencia de procesamiento
\end{itemize}
\end{minipage} \\
\hline
\textbf{FE-04: Cierre inesperado} & 
\begin{minipage}[t]{10.5cm}
\begin{itemize}
    \item Cierre forzado durante captura activa
    \item Sistema intenta guardar datos parciales automáticamente
    \item Registra estado de captura para recuperación
    \item En siguiente inicio, ofrece recuperar sesión interrumpida
\end{itemize}
\end{minipage} \\
\hline
\end{tabular}
}
\caption{Flujos de excepción del caso de uso CU-01}
\label{tab:cu01_flujos_excepcion}
\end{table}

\section{CU-02 - Registro temporal preciso durante experimentos}
\label{sec:cu02_registro_temporal}

\begin{table}[H]
\centering
\resizebox{\textwidth}{!}{
\rowcolors{2}{white}{udcgray!25}
\begin{tabular}{|p{3cm}|p{12cm}|}
\hline
\rowcolor{udcpink!25}
\textbf{Campo} & \textbf{Descripción} \\
\hline
\textbf{Identificador} & CU-02 \\
\hline
\textbf{Nombre} & Registro temporal preciso de acciones del usuario durante experimentos \\
\hline
\textbf{Actores} & Investigador neurocientífico (principal), Usuario del sistema (secundario) \\
\hline
\textbf{Descripción} & El sistema debe mantener precisión temporal de milisegundos en el registro de eventos para facilitar sincronización posterior con señales cerebrales \\
\hline
\textbf{Precondiciones} & 
\begin{minipage}[t]{11cm}
\begin{itemize}
    \item El sistema de captura está configurado y funcionando
    \item Se ha sincronizado el reloj del sistema
    \item El protocolo experimental está definido
    \item Los parámetros de precisión temporal están configurados
\end{itemize}
\end{minipage} \\
\hline
\textbf{Postcondiciones} & 
\begin{minipage}[t]{11cm}
\begin{itemize}
    \item Todos los eventos tienen timestamps de alta precisión
    \item Los datos están estructurados para análisis temporal
    \item Se mantiene trazabilidad de la secuencia temporal
    \item Los datos son exportables en formatos estándar
\end{itemize}
\end{minipage} \\
\hline
\end{tabular}
}
\caption{Especificación del caso de uso CU-02}
\label{tab:cu02_especificacion}
\end{table}

\subsection{Flujo principal}
\label{subsec:cu02_flujo_principal}

\begin{table}[H]
\centering
\resizebox{\textwidth}{!}{
\rowcolors{2}{white}{udcgray!25}
\begin{tabular}{|c|p{13cm}|}
\hline
\rowcolor{udcpink!25}
\textbf{Paso} & \textbf{Acción} \\
\hline
1 & El investigador neurocientífico configura parámetros de precisión temporal \\
\hline
2 & El sistema inicializa el reloj de alta precisión usando time.time() \\
\hline
3 & Se establecen marcadores temporales de inicio de experimento \\
\hline
4 & El investigador inicia el protocolo experimental \\
\hline
5 & El usuario del sistema comienza a interactuar según protocolo \\
\hline
6 & Cada evento se registra con timestamp preciso al momento de captura \\
\hline
7 & El sistema calcula y almacena intervalos entre eventos \\
\hline
8 & Se mantiene buffer circular para eventos de alta frecuencia \\
\hline
9 & Los timestamps se validan para detectar inconsistencias temporales \\
\hline
10 & El sistema genera metadatos de calidad temporal \\
\hline
11 & Al finalizar, se produce informe de precisión temporal \\
\hline
12 & Los datos se estructuran con timestamps como índice principal \\
\hline
\end{tabular}
}
\caption{Flujo principal del caso de uso CU-02}
\label{tab:cu02_flujo_principal}
\end{table}

\subsection{Flujos alternativos}
\label{subsec:cu02_flujos_alternativos}

\begin{table}[H]
\centering
\resizebox{\textwidth}{!}{
\rowcolors{2}{white}{udcgray!25}
\begin{tabular}{|p{3cm}|p{11cm}|}
\hline
\rowcolor{udcpink!25}
\textbf{Flujo} & \textbf{Descripción} \\
\hline
\textbf{FA-01: Calibración temporal} & 
\begin{minipage}[t]{10.5cm}
\begin{itemize}
    \item Antes del paso 4, se ejecuta calibración de latencia
    \item Sistema mide latencia propia mediante eventos sintéticos
    \item Se establece corrección temporal base
    \item Continúa con protocolo experimental ajustado
\end{itemize}
\end{minipage} \\
\hline
\textbf{FA-02: Marcadores experimentales} & 
\begin{minipage}[t]{10.5cm}
\begin{itemize}
    \item Durante el experimento, investigador inserta marcadores
    \item Marcadores incluyen timestamps de alta precisión
    \item Se etiquetan fases específicas del protocolo
    \item Facilita segmentación posterior de datos
\end{itemize}
\end{minipage} \\
\hline
\textbf{FA-03: Verificación tiempo real} & 
\begin{minipage}[t]{10.5cm}
\begin{itemize}
    \item Sistema monitoriza deriva temporal durante captura
    \item Detecta anomalías en secuencia temporal
    \item Alerta al investigador sobre problemas potenciales
    \item Mantiene log de calidad temporal
\end{itemize}
\end{minipage} \\
\hline
\end{tabular}
}
\caption{Flujos alternativos del caso de uso CU-02}
\label{tab:cu02_flujos_alternativos}
\end{table}

\subsection{Flujos de excepción}
\label{subsec:cu02_flujos_excepcion}

\begin{table}[H]
\centering
\resizebox{\textwidth}{!}{
\rowcolors{2}{white}{udcgray!25}
\begin{tabular}{|p{3cm}|p{11cm}|}
\hline
\rowcolor{udcpink!25}
\textbf{Excepción} & \textbf{Manejo} \\
\hline
\textbf{FE-01: Deriva temporal detectada} & 
\begin{minipage}[t]{10.5cm}
\begin{itemize}
    \item Sistema detecta inconsistencias en timestamps
    \item Registra evento de deriva con magnitud estimada
    \item Continúa captura aplicando corrección automática
    \item Marca datos afectados para revisión posterior
\end{itemize}
\end{minipage} \\
\hline
\textbf{FE-02: Sobrecarga temporal} & 
\begin{minipage}[t]{10.5cm}
\begin{itemize}
    \item Eventos llegan con frecuencia superior a capacidad procesamiento
    \item Sistema activa modo de preservación temporal
    \item Prioriza precisión sobre completitud de datos
    \item Documenta eventos descartados con timestamps
\end{itemize}
\end{minipage} \\
\hline
\textbf{FE-03: Fallo sincronización} & 
\begin{minipage}[t]{10.5cm}
\begin{itemize}
    \item Pérdida de referencia temporal del sistema
    \item Sistema continúa con timestamps relativos
    \item Registra punto de fallo para corrección posterior
    \item Mantiene coherencia interna de secuencia
\end{itemize}
\end{minipage} \\
\hline
\textbf{FE-04: Precisión degradada} & 
\begin{minipage}[t]{10.5cm}
\begin{itemize}
    \item Latencia del sistema supera umbrales aceptables
    \item Sistema alerta sobre degradación de precisión
    \item Ofrece opciones de optimización automática
    \item Documenta periodos de precisión reducida
\end{itemize}
\end{minipage} \\
\hline
\end{tabular}
}
\caption{Flujos de excepción del caso de uso CU-02}
\label{tab:cu02_flujos_excepcion}
\end{table}

\section{CU-04 - Herramientas básicas de visualización}
\label{sec:cu04_visualizacion_basica}

\begin{table}[H]
\centering
\resizebox{\textwidth}{!}{
\rowcolors{2}{white}{udcgray!25}
\begin{tabular}{|p{3cm}|p{12cm}|}
\hline
\rowcolor{udcpink!25}
\textbf{Campo} & \textbf{Descripción} \\
\hline
\textbf{Identificador} & CU-04 \\
\hline
\textbf{Nombre} & Herramientas básicas de visualización para investigación neurocientífica \\
\hline
\textbf{Actores} & Investigador neurocientífico (principal), Personal técnico de laboratorio (secundario) \\
\hline
\textbf{Descripción} & El sistema debe proporcionar visualizaciones en tiempo real que faciliten el monitoreo de actividad cerebral y eventos de entrada durante experimentos \\
\hline
\textbf{Precondiciones} & 
\begin{minipage}[t]{11cm}
\begin{itemize}
    \item El sistema está capturando datos de entrada o EEG
    \item La interfaz gráfica está inicializada correctamente
    \item Los visualizadores están configurados apropiadamente
    \item Hay datos disponibles para mostrar
\end{itemize}
\end{minipage} \\
\hline
\textbf{Postcondiciones} & 
\begin{minipage}[t]{11cm}
\begin{itemize}
    \item Las visualizaciones muestran datos actualizados en tiempo real
    \item Los investigadores pueden interpretar el estado del sistema
    \item Se mantiene historial visual configurable
    \item Las visualizaciones son exportables si es necesario
\end{itemize}
\end{minipage} \\
\hline
\end{tabular}
}
\caption{Especificación del caso de uso CU-04}
\label{tab:cu04_especificacion}
\end{table}

\subsection{Flujo principal}
\label{subsec:cu04_flujo_principal}

\begin{table}[H]
\centering
\resizebox{\textwidth}{!}{
\rowcolors{2}{white}{udcgray!25}
\begin{tabular}{|c|p{13cm}|}
\hline
\rowcolor{udcpink!25}
\textbf{Paso} & \textbf{Acción} \\
\hline
1 & El investigador inicia el sistema con módulos de visualización \\
\hline
2 & El sistema inicializa visualizadores según tipo de datos disponibles \\
\hline
3 & Se cargan configuraciones de visualización guardadas o por defecto \\
\hline
4 & Los visualizadores se conectan a flujos de datos correspondientes \\
\hline
5 & Comienza la actualización visual en tiempo real \\
\hline
6 & El investigador puede ajustar parámetros de visualización dinámicamente \\
\hline
7 & Los visualizadores adaptan su presentación según configuración \\
\hline
8 & Se mantiene historial visual según ventana temporal configurada \\
\hline
9 & El investigador monitoriza actividad y calidad de datos \\
\hline
10 & Sistema permite pausar/reanudar visualizaciones según necesidad \\
\hline
11 & Al finalizar, las configuraciones se guardan para sesiones futuras \\
\hline
\end{tabular}
}
\caption{Flujo principal del caso de uso CU-04}
\label{tab:cu04_flujo_principal}
\end{table}

\subsection{Flujos alternativos}
\label{subsec:cu04_flujos_alternativos}

\begin{table}[H]
\centering
\resizebox{\textwidth}{!}{
\rowcolors{2}{white}{udcgray!25}
\begin{tabular}{|p{3cm}|p{11cm}|}
\hline
\rowcolor{udcpink!25}
\textbf{Flujo} & \textbf{Descripción} \\
\hline
\textbf{FA-01: Visualización EEG multicanal} & 
\begin{minipage}[t]{10.5cm}
\begin{itemize}
    \item Sistema detecta datos EEG disponibles
    \item Inicializa visualizador multicanal con escala automática
    \item Muestra señales de hasta 16 canales simultáneamente
    \item Incluye métricas de calidad por canal
\end{itemize}
\end{minipage} \\
\hline
\textbf{FA-02: Visualización topográfica} & 
\begin{minipage}[t]{10.5cm}
\begin{itemize}
    \item Se activa mapa cerebral para distribución espacial
    \item Posiciona electrodos según sistema 10-20/10-10
    \item Codifica calidad de señal por colores
    \item Actualiza mapa según actividad en tiempo real
\end{itemize}
\end{minipage} \\
\hline
\textbf{FA-03: Visualización eventos entrada} & 
\begin{minipage}[t]{10.5cm}
\begin{itemize}
    \item Muestra teclas activas en representación visual
    \item Visualiza movimiento de ratón según modo configurado
    \item Actualiza estado de botones y dispositivos
    \item Mantiene historial de actividad reciente
\end{itemize}
\end{minipage} \\
\hline
\end{tabular}
}
\caption{Flujos alternativos del caso de uso CU-04}
\label{tab:cu04_flujos_alternativos}
\end{table}

\subsection{Flujos de excepción}
\label{subsec:cu04_flujos_excepcion}

\begin{table}[H]
\centering
\resizebox{\textwidth}{!}{
\rowcolors{2}{white}{udcgray!25}
\begin{tabular}{|p{3cm}|p{11cm}|}
\hline
\rowcolor{udcpink!25}
\textbf{Excepción} & \textbf{Manejo} \\
\hline
\textbf{FE-01: Sobrecarga visual} & 
\begin{minipage}[t]{10.5cm}
\begin{itemize}
    \item Datos llegan más rápido que capacidad de renderizado
    \item Sistema reduce frecuencia de actualización automáticamente
    \item Mantiene datos recientes priorizando fluidez visual
    \item Alerta sobre reducción de frecuencia visual
\end{itemize}
\end{minipage} \\
\hline
\textbf{FE-02: Fallo visualizador} & 
\begin{minipage}[t]{10.5cm}
\begin{itemize}
    \item Error en componente específico de visualización
    \item Sistema aísla visualizador problemático
    \item Continúa con visualizadores restantes funcionando
    \item Ofrece reinicialización del componente afectado
\end{itemize}
\end{minipage} \\
\hline
\textbf{FE-03: Pérdida datos visualización} & 
\begin{minipage}[t]{10.5cm}
\begin{itemize}
    \item Interrupción en flujo de datos hacia visualizadores
    \item Sistema muestra último estado conocido
    \item Indica visualmente pérdida de conectividad
    \item Restablece automáticamente al recuperar datos
\end{itemize}
\end{minipage} \\
\hline
\textbf{FE-04: Recursos insuficientes} & 
\begin{minipage}[t]{10.5cm}
\begin{itemize}
    \item Memoria o CPU insuficiente para visualización completa
    \item Sistema ofrece modos de visualización simplificados
    \item Permite desactivar visualizadores no críticos
    \item Monitoriza recursos y adapta automáticamente
\end{itemize}
\end{minipage} \\
\hline
\end{tabular}
}
\caption{Flujos de excepción del caso de uso CU-04}
\label{tab:cu04_flujos_excepcion}
\end{table}

\section{CU-05 - Configuración básica de dispositivos}
\label{sec:cu05_configuracion_dispositivos}

\begin{table}[H]
\centering
\resizebox{\textwidth}{!}{
\rowcolors{2}{white}{udcgray!25}
\begin{tabular}{|p{3cm}|p{12cm}|}
\hline
\rowcolor{udcpink!25}
\textbf{Campo} & \textbf{Descripción} \\
\hline
\textbf{Identificador} & CU-05 \\
\hline
\textbf{Nombre} & Configuración básica de dispositivos y validación de señal \\
\hline
\textbf{Actores} & Personal técnico de laboratorio (principal), Investigador neurocientífico (secundario) \\
\hline
\textbf{Descripción} & El sistema debe facilitar la configuración de dispositivos EEG y entrada, proporcionando validación de conectividad y calidad de señal \\
\hline
\textbf{Precondiciones} & 
\begin{minipage}[t]{11cm}
\begin{itemize}
    \item Los dispositivos están físicamente conectados
    \item Los drivers necesarios están instalados
    \item El personal técnico tiene conocimientos básicos del sistema
    \item BrainFlow está correctamente instalado para dispositivos EEG
\end{itemize}
\end{minipage} \\
\hline
\textbf{Postcondiciones} & 
\begin{minipage}[t]{11cm}
\begin{itemize}
    \item Los dispositivos están configurados correctamente
    \item La calidad de señal está validada
    \item Las configuraciones están guardadas para uso futuro
    \item El sistema está listo para experimentos
\end{itemize}
\end{minipage} \\
\hline
\end{tabular}
}
\caption{Especificación del caso de uso CU-05}
\label{tab:cu05_especificacion}
\end{table}

\subsection{Flujo principal}
\label{subsec:cu05_flujo_principal}

\begin{table}[H]
\centering
\resizebox{\textwidth}{!}{
\rowcolors{2}{white}{udcgray!25}
\begin{tabular}{|c|p{13cm}|}
\hline
\rowcolor{udcpink!25}
\textbf{Paso} & \textbf{Acción} \\
\hline
1 & El personal técnico accede al módulo de configuración \\
\hline
2 & El sistema escanea dispositivos disponibles automáticamente \\
\hline
3 & Se presenta lista de dispositivos detectados y compatibles \\
\hline
4 & El técnico selecciona dispositivo EEG específico a configurar \\
\hline
5 & Sistema inicializa conexión con parámetros predeterminados \\
\hline
6 & Se ejecuta test básico de conectividad \\
\hline
7 & El técnico configura parámetros específicos (frecuencia muestreo, canales) \\
\hline
8 & Sistema valida configuración y establece conexión estable \\
\hline
9 & Se ejecuta validación de calidad de señal por canal \\
\hline
10 & El técnico revisa métricas de calidad y ajusta si necesario \\
\hline
11 & Configuración validada se guarda como perfil de dispositivo \\
\hline
12 & Sistema confirma configuración exitosa y dispositivo listo \\
\hline
\end{tabular}
}
\caption{Flujo principal del caso de uso CU-05}
\label{tab:cu05_flujo_principal}
\end{table}

\subsection{Flujos alternativos}
\label{subsec:cu05_flujos_alternativos}

\begin{table}[H]
\centering
\resizebox{\textwidth}{!}{
\rowcolors{2}{white}{udcgray!25}
\begin{tabular}{|p{3cm}|p{11cm}|}
\hline
\rowcolor{udcpink!25}
\textbf{Flujo} & \textbf{Descripción} \\
\hline
\textbf{FA-01: Configuración placa sintética} & 
\begin{minipage}[t]{10.5cm}
\begin{itemize}
    \item En paso 4, técnico selecciona placa sintética BrainFlow
    \item Sistema configura generación de datos sintéticos
    \item Permite configurar patrones de señal para pruebas
    \item Útil para validación sin hardware real
\end{itemize}
\end{minipage} \\
\hline
\textbf{FA-02: Configuración dispositivos entrada} & 
\begin{minipage}[t]{10.5cm}
\begin{itemize}
    \item Configuración paralela de teclado y ratón
    \item Validación de permisos de acceso del sistema
    \item Configuración de filtros de teclas relevantes
    \item Test de responsividad de dispositivos
\end{itemize}
\end{minipage} \\
\hline
\textbf{FA-03: Múltiples perfiles} & 
\begin{minipage}[t]{10.5cm}
\begin{itemize}
    \item Creación de múltiples perfiles para mismo dispositivo
    \item Perfiles específicos para diferentes protocolos experimentales
    \item Cambio rápido entre configuraciones guardadas
    \item Exportación/importación de perfiles entre sistemas
\end{itemize}
\end{minipage} \\
\hline
\end{tabular}
}
\caption{Flujos alternativos del caso de uso CU-05}
\label{tab:cu05_flujos_alternativos}
\end{table}

\subsection{Flujos de excepción}
\label{subsec:cu05_flujos_excepcion}

\begin{table}[H]
\centering
\resizebox{\textwidth}{!}{
\rowcolors{2}{white}{udcgray!25}
\begin{tabular}{|p{3cm}|p{11cm}|}
\hline
\rowcolor{udcpink!25}
\textbf{Excepción} & \textbf{Manejo} \\
\hline
\textbf{FE-01: Dispositivo no detectado} & 
\begin{minipage}[t]{10.5cm}
\begin{itemize}
    \item Sistema no encuentra dispositivos compatibles
    \item Muestra guía de resolución de problemas específica
    \item Verifica drivers y configuración de puertos
    \item Ofrece modo manual de configuración
\end{itemize}
\end{minipage} \\
\hline
\textbf{FE-02: Fallo conectividad} & 
\begin{minipage}[t]{10.5cm}
\begin{itemize}
    \item Error al establecer conexión con dispositivo
    \item Sistema intenta reconexión automática con parámetros alternativos
    \item Proporciona diagnóstico detallado del fallo
    \item Sugiere verificaciones físicas de conexión
\end{itemize}
\end{minipage} \\
\hline
\textbf{FE-03: Calidad señal deficiente} & 
\begin{minipage}[t]{10.5cm}
\begin{itemize}
    \item Métricas de calidad por debajo de umbrales aceptables
    \item Sistema identifica canales específicos problemáticos
    \item Proporciona guía de mejora de impedancia/conectividad
    \item Permite ajuste de umbrales según contexto experimental
\end{itemize}
\end{minipage} \\
\hline
\textbf{FE-04: Conflicto configuración} & 
\begin{minipage}[t]{10.5cm}
\begin{itemize}
    \item Parámetros configurados no compatibles con dispositivo
    \item Sistema sugiere configuraciones alternativas válidas
    \item Permite configuración asistida basada en capacidades detectadas
    \item Valida configuración antes de aplicar cambios
\end{itemize}
\end{minipage} \\
\hline
\end{tabular}
}
\caption{Flujos de excepción del caso de uso CU-05}
\label{tab:cu05_flujos_excepcion}
\end{table}